\documentclass[11pt]{article}
\usepackage{fontspec,amsmath,amssymb,amsthm,hyperref}
\usepackage[margin=1in]{geometry}
\setlength{\emergencystretch}{3em}
\title{Actual kinetic flows under measure-preserving source changes}
\author{LangevinLean mathematical exposition}
\date{}
\begin{document}\maketitle
\noindent Source module: \texttt{KineticLangevin/FlowPullback.lean}. Compile this proof document with XeLaTeX or LuaLaTeX. The accompanying HTML pairs the mathematical argument with the exact Lean statements and records the scope of the formalization.


Changing a probability source by a measure-preserving projection preserves the actual random variables and their output laws. This module proves that the genuine lifted nonlinear drift and the selected kinetic integral solution commute with that change of source. It supplies coupling-independent continuous-flow laws without assuming samplewise continuity of an arbitrary L\ensuremath{{}^{2}} representative.


\section{Define continuous linear source and phase pullbacks}


Measure-preserving composition is a linear isometry on L\ensuremath{{}^{2}}. The source pullback is this isometry viewed as a continuous linear map, and the phase pullback applies it separately to velocity and position. Their representative theorems identify the resulting random variables with actual composition. Pushforward composition proves preservation of every output law. A further representative argument proves that two successive pullbacks equal pullback along the composed projection.



\noindent Lean declarations: \texttt{sourcePullback}, \texttt{phasePullback}, \texttt{sourcePullback\_coe}, \texttt{phasePullback\_coe}, \texttt{sourcePullback\_map}, \texttt{sourcePullback\_comp}.

\section{Commute the actual nonlinear drift with pullback}


The lifted vector field is pointwise application of the original Lipschitz field. Lean compares its actual representatives before and after source projection, obtaining exact equality in L\ensuremath{{}^{2}}. Combining this with linearity proves the kinetic drift commutes with the phase pullback, including both friction and force terms and the position equation.



\noindent Lean declarations: \texttt{lpDrift\_sourcePullback}, \texttt{kineticDrift\_phasePullback}.

\section{Identify the selected actual integral solutions}


Apply the phase pullback to the original kinetic integral equation. Continuity and linearity commute it with the Bochner time integral, and the drift commutation identity produces the target-space kinetic equation. The transformed initial condition and additive input are exact. The proved integral-solution uniqueness theorem therefore identifies the selected target-space solution with the pulled original solution at every time in the horizon. This equality transfers laws between an independent initial-state/noise product and any larger initial-coupling extension.



The reusable integral-equation theorem proves this transformation for any genuine continuous kinetic solution before specializing to selected flows.



\noindent Lean declarations: \texttt{kinetic\_integral\_sourcePullback}, \texttt{kineticSegmentFlow\_sourcePullback}.

\end{document}
